\newpage
\section{Introduction}
\label{sec:intro}

Chatbot assistants powered by \emph{large language models} (LLMs) have
transformed how users access and interact with information. Since the early
2020s, adoption has grown rapidly. By September 2025, one of the largest 
foundation labs reported over 700 million weekly active users generating billions 
of messages per week. Usage is particularly concentrated among 18--25-year-olds, 
who account for roughly 46\% of all messages~\cite{chatgptusage2025}. This youth 
adoption suggests a structural shift in how knowledge is queried and sought:
natural language interaction increasingly complements traditional search
interfaces such as folder hierarchies or keyword-based matching.

Most widely used LLM assistants are cloud-hosted and operated by third
parties. While convenient, this architecture raises important concerns.
Sensitive organisational data may leave controlled environments and the
cost of access depends on external providers whose pricing and availability
are not guaranteed. 

For many organisations, these considerations make local deployment
attractive. Running language models on controlled infrastructure preserves
data ownership and reduces dependence on external services.

However, a language model alone is insufficient for many local applications.
Its knowledge is largely fixed at training time and does not
include private documents, internal procedures or other organisation-specific
information. Conversely, a document database alone does not provide an
interactive interface for exploring knowledge.

\emph{Retrieval-Augmented Generation} (RAG) systems address this limitation
by combining both components: given a user query, the system retrieves
relevant information from internal sources and provides it as context for
the model’s response.

\begin{figure}[h!]
    \centering
        \centering
        \includegraphics[width=0.5\textwidth]{rag-principle.png}
        \caption{Retrieval-augmented generation overview.}
        \label{fig:rag-principle}
\end{figure}

\subsection{Motivation}

Deploying language models locally also changes the underlying cost
structure. Instead of per-call access to cloud APIs,
organisations must operate their own computational infrastructure.

While training large models requires substantial computational resources,
it occurs relatively infrequently. In contrast, inference---the repeated
execution of the model for user queries---occurs continuously and therefore
dominates operational cost in many deployments. Industry analyses confirm
the growth: as training capabilities mature, inference will represent the
largest share of compute demand and operational expenditure in deployed
AI systems~\cite{inferenceshift2025, morganstanleyinference2025}.

Advanced RAG pipelines often include optional modules intended to improve
retrieval or answer quality. These may involve query expansion,
hybrid search or reranking components. In practice, such modules are
frequently enabled in static configurations where every query triggers
the full pipeline regardless of its complexity.

This design is simple and performance-oriented, but implicitly assumes that
additional computation is always worth the cost. In local deployments with
fixed hardware resources, this assumption can become inefficient.

Despite the growing deployment of RAG systems, their operational efficiency,
particularly the energy cost of modular pipelines, remains relatively
underexplored.

\subsection{Problem Statement}

This thesis studies RAG pipelines from a system perspective with energy as
a first-class metric. The central hypothesis is that \textbf{not all queries
require the same RAG modules} and that \textbf{selectively activating
computationally expensive modules can reduce energy consumption while
preserving acceptable answer quality}.

A deeper question underlies this hypothesis: can a router decide which
modules to activate \emph{before} observing their output?  If the
information needed to route is produced by the very computation that
routing aims to skip, pre-execution routing faces a structural ceiling.
We investigate whether this circular dependency holds for module-level
activation in RAG systems.

\subsection{Research Objectives}

To test these hypotheses, we
(i) design a measurement framework for distributed RAG energy attribution,
(ii) quantify the energy--quality trade-offs of optional modules under
controlled configurations,
(iii) formalize per-query configuration control as an optimization problem
and (iv) systematically evaluate routing strategies to determine the
achievable energy savings frontier.

\subsection{Contributions}

The main contributions of this thesis are:

\begin{itemize}
    \item \textbf{Contribution 1}: An energy measurement framework for distributed RAG
          systems that attributes cost per module, per node and per pipeline
          stage, reusable for modular RAG deployments.

    \item \textbf{Contribution 2}: Empirical evidence that always-on module activation is
          inefficient, with energy costs spanning 5.6$\times$ for quality
          differences below 0.06 points across eight configurations.

    \item \textbf{Contribution 3}: A systematic routing investigation across eight
          experiments (five strategies and three explorations) that
          quantifies the signal boundary for pre-execution module
          prediction, showing that query-only features provide weak
          discriminative power for module-level configuration control.

    \item \textbf{Contribution 4}: Characterisation of the oracle routing ceiling (75\%
          energy savings in the full configuration space, reduced to 19\%
          in the non-HyDE space tested for routing) and quantification of
          the information gap between oracle (100\% config accuracy) and
          achievable pre-execution routing (39.5\%).
\end{itemize}

Contributions~1 and~2 provide the measurement infrastructure and empirical
evidence.  Contributions~3 and~4 constitute the central finding: a
characterised structural ceiling for pre-execution module routing, explained
by a circular dependency between routing decisions and pipeline outputs.

